Algorithmic redistricting --- the use of computational optimization to draw electoral district boundaries --- has emerged as a promising alternative to legislative gerrymandering \citep{dellib2026edgeweighted, cirincione2000measuring, cho2017measuring}. However, the existing literature is almost entirely focused on the United States, where single-member districts elected by plurality voting create specific incentives for partisan manipulation. This geographic and institutional parochialism raises a fundamental question: is algorithmic redistricting a solution to a uniquely American problem, or does it offer benefits transferable to other democratic systems?

The question matters for several reasons. First, many established democracies use independent boundary commissions (UK, Canada, Australia, New Zealand) that already pursue the same objectives as algorithmic approaches --- compactness, population equality, and community preservation --- but through different processes. Comparing algorithmic outcomes to commission-drawn boundaries provides a direct test of algorithmic redistricting's incremental value. Second, the rise of multi-member proportional systems in established democracies (and proposals to adopt them in others) requires redistricting methodology that accounts for seat magnitude. Ireland's 43 constituencies return 3--5 seats each; Malta's 13 constituencies each return 5 seats. The balance criterion, compactness metric, and proportionality measurement all differ substantively from single-member equivalents.

Third, the RPLAN open interchange format \citep{redistcli2026} enables, for the first time, a standardized cross-national redistricting dataset. By defining constituency assignments as 11-character geographic unit identifiers mapped to district numbers --- independent of the electoral system, seat magnitude, or administrative geography --- RPLAN allows the same analysis pipeline to operate across countries using Census-equivalent data.

We make three contributions. First, we demonstrate that edge-weighted recursive bisection achieves compactness improvements of 10--22\% over enacted boundaries across all six countries studied, suggesting the algorithm's benefits are not contingent on American-specific institutional features. Second, we introduce per-seat balance and D'Hondt-based proportionality analysis for multi-member systems, demonstrating that geographic optimization and electoral proportionality are compatible objectives. Third, we identify which institutional features correlate with larger algorithmic improvements: countries with less independent redistricting commissions (Canada, Germany) show larger gains than those with stronger commissions (UK, New Zealand), consistent with the hypothesis that algorithmic redistricting is most valuable where political influence on boundary-drawing is greatest.

The paper proceeds as follows. Section 2 describes the comparative redistricting institutions and data for each country. Section 3 presents the algorithmic methodology, including extensions to multi-member systems. Section 4 reports compactness, balance, and proportionality results. Section 5 discusses policy implications and cross-national lessons.
